% !TeX root = ../navier-stokes-manuscript.tex
% ============================================================
% 1.4 Silver Standard Proof
% ============================================================
\subsection{Silver Standard Proof}\label{sub:SilverStandardProof}

Looking for the missing estimates to satisfy the Gold-standard requirements is really difficult.
No one's been able to do it in almost a hundred years, and once you dig your teeth into the problem, it becomes tempting to find any other way of proving smoothness.

Assume you already solved the problem using the Gold standard method.

That would mean smoothness itself is a property of the Navier-Stokes equations.
If smoothness is a property of the fluid itself, then it must've come from the only other properties in the fluid; velocity, pressure, viscosity, incompressibility, and non-linear transport, all acting in unison to keep the fluid smooth.
Which means any hypothetical “blow-up” scenario (or any case that breaks smoothness) should not be reachable from those equations, and any imagined blow-up scenario must not be describing the same fluid from the original problem statement.
If every possible loss of smoothness stops satisfying the original problem, or stops being about the same fluid, then none of those scenarios are valid Navier-Stokes counterexamples.
If there are no valid counterexamples, then the fluid must be smooth.

\[
  \begin{gathered}
  \text{If Navier--Stokes}\Longrightarrow\text{smooth},
  \\
  \text{then not smooth}\Longrightarrow\text{not Navier--Stokes}.
  \end{gathered}
\]

The Silver Standard proof is a proof by contrapositive.
Proving that every non-smooth outcome cannot be reached from the original Navier-Stokes equations, proves that the original equations always produce a smooth fluid.
The beauty of the Silver proof method is that it no longer asks \textit{“how is the fluid smooth?”}, it only asks \textit{“is it smooth?”}; which reveals why the Gold proof method is so hard — Gold requires an explanation of mechanism (and hence the reason for the podium-titles).
So the upside here, for Silver, is that the burden now shifts away from trying to find missing estimates.
The downside is exhaustively handling every non-smooth outcome.
Coming up with a principle that handles every case is the key to simplifying that burden.
